\documentclass[11pt]{article}
\usepackage{fontspec,amsmath,amssymb,amsthm,hyperref}
\usepackage[margin=1in]{geometry}
\setlength{\emergencystretch}{3em}
\title{Exact iid variance and the empirical parameter force}
\author{LangevinLean mathematical exposition}
\date{}
\begin{document}\maketitle
\noindent Source module: \texttt{ProximalLangevin/MeanFieldVariance.lean}. Compile this proof document with XeLaTeX or LuaLaTeX. The accompanying HTML pairs the mathematical argument with the exact Lean statements and records the scope of the formalization.


On a probability space, square-integrable Hilbert-valued variables have integrable inner products. On the actual finite product law, independent centered samples yield an exact \(1/N\) variance identity.


\section{Cancel every off-diagonal covariance}


Realize the supplied \(L^2\) functions as \(L^2\) elements; Cauchy–Schwarz gives integrable pairwise inner products. Expand

\[\mathbb E\left\|\sum_iX_i\right\|^2=\sum_{i,j}\mathbb E\langle X_i,X_j\rangle.\]

For \(i\ne j\), independence factors the expectation into \(\langle\mathbb E X_i,\mathbb E X_j\rangle=0\). The diagonal terms are \(\mathbb E\|X_i\|^2\). Finite sums and established integrability justify the expansion.




Expand the Hilbert square as \(\|\sum_iX_i\|^2=\sum_{i,j}\langle X_i,X_j\rangle\). Each mixed product is integrable by the \(L^2\) inner-product bound, which justifies moving the finite double sum through the integral. For \(i\ne j\), independence factors \(\mathbb E\langle X_i,X_j\rangle\) into \(\langle\mathbb EX_i,\mathbb EX_j\rangle=0\). For \(i=j\), it is \(\mathbb E\|X_i\|^2\). The local facts \texttt{hint}, \texttt{hcross} and \texttt{hexpand} establish these prerequisites and identities separately. Centering without independence would leave the off-diagonal covariances.




\textbf{Correspondence with the displayed declarations.}


\begin{itemize}

\item \texttt{integrable\_inner\_of\_memLp}: Every inner product of two L\ensuremath{{}^{2}} random variables is integrable.
\item \texttt{integral\_norm\_sum\_sq\_centered\_independent}: Centered independent Hilbert-valued samples have exactly additive second moments. Independence is used for each mixed term.
\end{itemize}


\noindent Lean declarations: \texttt{integrable\_inner\_of\_memLp}, \texttt{integral\_norm\_sum\_sq\_centered\_independent}.

\section{Use the actual finite product for an exact sample-mean variance}


For a square-integrable observable \(\phi\) and a probability law \(\mu\), the finite product measure \(\mu^{\otimes N}\) constructs the independent samples. The proof centers \(\phi\), expands the squared Hilbert norm, and uses independence and zero mean to cancel every off-diagonal inner product. The exact result is


\[
\mathbb E_{\mu^{\otimes N}}
 \left\|\frac1N\sum_{i=1}^N\phi(X_i)-\mathbb E_\mu\phi\right\|^2
 =\frac1N\mathbb E_\mu\|\phi-\mathbb E_\mu\phi\|^2.
\]


Taking square roots proves the \(N^{-1/2}\) sampling factor. The actual parameter-gradient observable is square integrable under finite-P\ensuremath{{}_{2}} latent laws, so the formula applies to the mean-field force.



The source’s exchangeability premise alone does not cancel these covariances. A quantitative chaos estimate needs controlled initialization, such as the actual iid construction used here. This exact sampling identity is a proved ingredient; by itself it is not a propagation-of-chaos theorem.




On the explicit product probability \(P=\mu^{\otimes N}\), set \(C(a)=\varphi(a)-\mathbb E_\mu\varphi\) and \(X_i(\omega)=C(\omega_i)\). Measure-preserving coordinate evaluation proves that each \(X_i\) has the same finite second moment, zero mean and law, and the product construction proves their joint independence. The exact centered-mean identity is \[\frac1N\sum_i\varphi(\omega_i)-\mathbb E\varphi=\frac1N\sum_iX_i(\omega).\] Thus the previous variance identity gives \[\mathbb E_P\left\|\frac1N\sum_i\varphi(\omega_i)-\mathbb E\varphi\right\|^2=\frac{\mathbb E_\mu\|\varphi-\mathbb E\varphi\|^2}{N}.\] In the code, \texttt{hXi}, \texttt{hind}, \texttt{hmean}, \texttt{hvar} and \texttt{hid} verify each of these assertions. Taking the nonnegative square root gives the RMS version.




\textbf{Correspondence with the displayed declarations.}


\begin{itemize}

\item \texttt{empiricalAverage}: The empirical mean on the explicitly specified finite sample space.
\item \texttt{iid\_empirical\_variance}: The actual iid finite product gives the exact centered empirical variance. All moments are Bochner integrals, including the mean used for centering.
\item \texttt{iid\_empirical\_rms}: The same exact identity written as an RMS sampling error.
\end{itemize}


\noindent Lean declarations: \texttt{empiricalAverage}, \texttt{iid\_empirical\_variance}, \texttt{iid\_empirical\_rms}.

\section{Specialize the variance to a true gradient observable}


Fix \(\theta\) and let \(\phi(x)=\nabla_\theta U(\theta,x)\). A Lipschitz joint gradient makes this observable Lipschitz, hence square integrable under any finite-second-moment latent law \(\mu\). Its expectation is exactly \(b_U(\theta,\mu)\). Substitute into the iid formula to obtain

\[\mathbb E\left\|\frac1N\sum_i\nabla_\theta U(\theta,X_i)-b_U(\theta,\mu)\right\|^2=\frac1N\mathbb E\|\nabla_\theta U(\theta,X)-b_U(\theta,\mu)\|^2.\]


For \(\varphi(x)=\nabla_\theta U(\theta,x)\), the parameter projection is nonexpansive and the full score is \(L\)-Lipschitz, so \(\varphi\) is \(L\)-Lipschitz in the latent variable. A latent law with finite second moment therefore makes this observable square integrable. Its Bochner mean is exactly \texttt{meanFieldDrift}, so the general iid identity applies without assuming a separate variance bound for this particular gradient observable.




\textbf{Correspondence with the displayed declarations.}


\begin{itemize}

\item \texttt{iid\_parameterGradient\_variance}: Exact iid empirical parameter-gradient variance for the actual joint potential. The copies are independent under the displayed finite product.
\end{itemize}


\noindent Lean declarations: \texttt{iid\_parameterGradient\_variance}.

\end{document}
